\chapter{การออกแบบพฤติกรรมระบบ (Behavioral System Design - UML Diagrams)}

บทนี้จะอธิบายรายละเอียดการออกแบบพฤติกรรมของระบบบริหารจัดการยิม (GymBro Management System) โดยใช้แผนภาพ Unified Modeling Language (UML) แบบจำลองพฤติกรรมแสดงให้เห็นว่าระบบมีปฏิสัมพันธ์กับหน่วยงานภายนอก (ผู้ใช้) อย่างไร และส่วนประกอบภายในทำงานร่วมกันอย่างไรเพื่อให้เป็นไปตามข้อกำหนดการทำงาน ส่วนนี้ครอบคลุมแผนภาพยูสเคส (Use Case Diagrams), แผนภาพกิจกรรม (Activity Diagrams) และแผนภาพลำดับ (Sequence Diagrams)

\section{แผนภาพยูสเคสและคำอธิบายอย่างละเอียด (Use Case Diagram \& Detailed Descriptions)}
แผนภาพยูสเคสให้มุมมองระดับสูงของการทำงานของระบบจากมุมมองของผู้กระทำภายนอก (External Actors) โดยกำหนดขอบเขตของระบบและการปฏิสัมพันธ์ระหว่างผู้กระทำและยูสเคส

\subsection{ผู้กระทำ (Actors)}
ระบบเกี่ยวข้องกับผู้กระทำหลักสองกลุ่ม:
\begin{itemize}
    \item \textbf{ผู้ดูแลระบบ (Administrator)}: ผู้ใช้ที่มีสิทธิพิเศษ รับผิดชอบในการจัดการข้อมูลระบบ รวมถึงระเบียนสมาชิก โปรไฟล์เทรนเนอร์ สถานะอุปกรณ์ และตรวจสอบบันทึกระบบ
    \item \textbf{สมาชิก (Member)}: ผู้ใช้ที่ลงทะเบียนแล้ว ซึ่งโต้ตอบกับระบบเพื่อดูตารางเวลา จองเซสชันการฝึกซ้อม และจัดการโปรไฟล์ส่วนตัว
\end{itemize}

\subsection{รายการยูสเคส (Use Case List)}
\textbf{ยูสเคสของผู้ดูแลระบบ:}
\begin{itemize}
    \item เข้าสู่ระบบ (Admin Authentication)
    \item จัดการสมาชิก (ดู, แก้ไข, ลบ)
    \item จัดการเทรนเนอร์ (เพิ่ม, แก้ไข, ลบ)
    \item จัดการอุปกรณ์ (เพิ่ม, แก้ไขสถานะ, ลบ)
    \item ดูแดชบอร์ดและสถิติระบบ
    \item ดูบันทึกระบบ (System Logs)
    \item จัดการเซสชันการจองทั้งหมด (ดู, ยกเลิก)
\end{itemize}

\textbf{ยูสเคสของสมาชิก:}
\begin{itemize}
    \item ลงทะเบียน (สร้างบัญชี)
    \item เข้าสู่ระบบ (Member Authentication)
    \item ดูแดชบอร์ดส่วนตัวและตารางเวลา
    \item จองเซสชันการฝึกซ้อม
    \item ยกเลิกการจองของตนเอง
    \item ดูโปรไฟล์ผู้ใช้
\end{itemize}

\subsection{คำอธิบายยูสเคสอย่างละเอียด (Detailed Use Case Descriptions)}
ตารางต่อไปนี้อธิบายกระบวนการสำคัญ 4 ประการอย่างละเอียด

% Table 1: User Authentication
\begin{longtable}{|p{3cm}|p{12cm}|}
\hline
\textbf{รหัสยูสเคส} & UC-01 \\ \hline
\textbf{ชื่อยูสเคส} & \textbf{การยืนยันตัวตนผู้ใช้ (Register/Login)} \\ \hline
\textbf{ผู้กระทำหลัก} & สมาชิก, ผู้ดูแลระบบ \\ \hline
\textbf{เงื่อนไขก่อนเริ่ม} & ผู้ใช้เข้าถึงเว็บแอปพลิเคชัน GymBro \\ \hline
\textbf{สถานการณ์หลัก (Success Scenario)} & 
\begin{enumerate}
    \item ผู้ใช้ไปยังหน้าเข้าสู่ระบบ (Login)
    \item ผู้ใช้กรอกอีเมลและเบอร์โทรศัพท์ที่ลงทะเบียนไว้
    \item ผู้ใช้กดปุ่ม "Sign In"
    \item ระบบตรวจสอบข้อมูลกับฐานข้อมูล
    \item ระบบดึงบทบาทของผู้ใช้ (Admin หรือ Member)
    \item ระบบสร้าง Session Token (เก็บใน localStorage)
    \item ระบบนำทางผู้ใช้ไปยัง Dashboard ที่เหมาะสม (Admin Dashboard หรือ Member Schedule)
\end{enumerate} \\ \hline
\textbf{ทางเลือกอื่น (Alternative Flows)} & 
\textbf{A1: ข้อมูลไม่ถูกต้อง}
\begin{enumerate}
    \item ระบบไม่พบข้อมูลที่ตรงกัน
    \item ระบบแสดงข้อความแจ้งเตือน "Invalid credentials"
    \item ผู้ใช้ยังคงอยู่ที่หน้า Login เพื่อลองใหม่
\end{enumerate}
\textbf{A2: การลงทะเบียนผู้ใช้ใหม่}
\begin{enumerate}
    \item ผู้ใช้คลิก "Register here"
    \item ผู้ใช้กรอกชื่อ-นามสกุล, อีเมล, และเบอร์โทรศัพท์
    \item ระบบตรวจสอบอีเมลซ้ำ
    \item ระบบสร้างระเบียนสมาชิกใหม่ด้วยสถานะ "Pending" (หรือ Active ตามค่าเริ่มต้น)
    \item ระบบนำทางผู้ใช้ไปหน้า Login
\end{enumerate} \\ \hline
\textbf{เงื่อนไขหลังจบ} & ผู้ใช้ได้รับการยืนยันตัวตนและเข้าถึงฟีเจอร์ตามสิทธิ์ \\ \hline
\end{longtable}

% Table 2: Equipment/Class Booking
\begin{longtable}{|p{3cm}|p{12cm}|}
\hline
\textbf{รหัสยูสเคส} & UC-04 \\ \hline
\textbf{ชื่อยูสเคส} & \textbf{การจองอุปกรณ์/คลาสเรียน (Equipment/Class Booking)} \\ \hline
\textbf{ผู้กระทำหลัก} & สมาชิก \\ \hline
\textbf{เงื่อนไขก่อนเริ่ม} & สมาชิกเข้าสู่ระบบแล้วและอยู่ที่หน้า User Dashboard \\ \hline
\textbf{สถานการณ์หลัก (Success Scenario)} & 
\begin{enumerate}
    \item สมาชิกคลิก "Book Session"
    \item ระบบแสดงหน้าต่างการจอง (Booking Modal)
    \item สมาชิกเลือกวันที่, เวลาเริ่ม, และระยะเวลา
    \item ระบบเรียก \textbf{อัลกอริทึมตรวจสอบเวลาซ้อนทับ} (UC-05)
    \item ระบบอัปเดตรายการเทรนเนอร์และอุปกรณ์:
    \begin{itemize}
        \item รายการที่ว่างจะเลือกได้
        \item รายการที่ไม่ว่างหรือซ่อมบำรุงจะถูกปิดการใช้งาน (Disabled)
    \end{itemize}
    \item สมาชิกเลือกเทรนเนอร์และอุปกรณ์ที่ว่าง
    \item สมาชิกคลิก "Confirm Booking"
    \item ระบบบันทึกการจองลงฐานข้อมูล
    \item ระบบแสดงแจ้งเตือนสำเร็จและรีเฟรชตารางเวลา
\end{enumerate} \\ \hline
\textbf{ทางเลือกอื่น (Alternative Flows)} & 
\textbf{A1: ไม่มีความพร้อมให้บริการ}
\begin{enumerate}
    \item สมาชิกเลือกเวลาที่เทรนเนอร์ทุกคนไม่ว่าง
    \item ระบบปิดการใช้งานตัวเลือกเทรนเนอร์ทั้งหมด
    \item สมาชิกต้องเลือกช่วงเวลาอื่น
\end{enumerate} \\ \hline
\textbf{เงื่อนไขหลังจบ} & ระเบียน `TrainingSession` ใหม่ถูกสร้างขึ้น เชื่อมโยงสมาชิก เทรนเนอร์ และอุปกรณ์ \\ \hline
\end{longtable}

% Table 3: Time Conflict Detection
\begin{longtable}{|p{3cm}|p{12cm}|}
\hline
\textbf{รหัสยูสเคส} & UC-05 \\ \hline
\textbf{ชื่อยูสเคส} & \textbf{อัลกอริทึมตรวจสอบเวลาซ้อนทับ (Time Conflict Detection)} \\ \hline
\textbf{ผู้กระทำหลัก} & ระบบ (กระบวนการภายใน) \\ \hline
\textbf{เงื่อนไขก่อนเริ่ม} & ผู้ใช้เลือกช่วงเวลาที่ต้องการจอง \\ \hline
\textbf{สถานการณ์หลัก (Success Scenario)} & 
\begin{enumerate}
    \item ระบบคำนวณ `startTime` และ `endTime` ที่ร้องขอ
    \item ระบบค้นหาฐานข้อมูลสำหรับเซสชันที่มีอยู่ทั้งหมดที่ทับซ้อนกับช่วงเวลานี้
    \item ตรรกะ: `(ExistingStart < NewEnd) AND (ExistingEnd > NewStart)`
    \item ระบบระบุว่าเทรนเนอร์และอุปกรณ์ใดเกี่ยวข้องกับเซสชันที่ทับซ้อน
    \item ระบบส่งคืนรายการ ID ที่ "ไม่ว่าง" ไปยัง Frontend
    \item Frontend ปิดการใช้งาน ID เหล่านี้ใน UI การเลือก
\end{enumerate} \\ \hline
\textbf{ทางเลือกอื่น (Alternative Flows)} & N/A (เป็นกระบวนการตรวจสอบเบื้องหลัง) \\ \hline
\textbf{เงื่อนไขหลังจบ} & ผู้ใช้ถูกป้องกันจากการจองทรัพยากรซ้อนทับ \\ \hline
\end{longtable}

% Table 4: Admin System Logging
\begin{longtable}{|p{3cm}|p{12cm}|}
\hline
\textbf{รหัสยูสเคส} & UC-08 \\ \hline
\textbf{ชื่อยูสเคส} & \textbf{การจัดการข้อมูลและการบันทึกระบบของผู้ดูแล (System Logging)} \\ \hline
\textbf{ผู้กระทำหลัก} & ผู้ดูแลระบบ \\ \hline
\textbf{เงื่อนไขก่อนเริ่ม} & Admin เข้าสู่ระบบแล้ว \\ \hline
\textbf{สถานการณ์หลัก (Success Scenario)} & 
\begin{enumerate}
    \item Admin ทำการดำเนินการสำคัญ (เช่น ลบสมาชิก, แก้ไขเทรนเนอร์)
    \item ระบบดำเนินการกับฐานข้อมูล
    \item เมื่อสำเร็จ ระบบสร้างรายการ Log ประกอบด้วย: Timestamp, Action Type, Admin ID, และ Details
    \item ระบบเพิ่มรายการนี้ลงในไฟล์ `logs.json` หรือตาราง Log ในฐานข้อมูล
    \item Admin ไปยังหน้า "System Logs"
    \item ระบบอ่านและแสดงประวัติการใช้งาน
\end{enumerate} \\ \hline
\textbf{ทางเลือกอื่น (Alternative Flows)} & 
\textbf{A1: ข้อผิดพลาดไฟล์ Log}
\begin{itemize}
    \item หากระบบไม่สามารถเขียนไฟล์ Log ได้ จะแสดงข้อผิดพลาดที่ Server Console แต่ยังคงดำเนินการตามคำสั่งผู้ใช้เพื่อให้งานไม่สะดุด
\end{itemize} \\ \hline
\textbf{เงื่อนไขหลังจบ} & บันทึกการกระทำของผู้ดูแลระบบถูกเก็บรักษาไว้เพื่อการตรวจสอบ \\ \hline
\end{longtable}

\newpage
\section{แผนภาพกิจกรรม (Activity Diagrams)}
แผนภาพกิจกรรมอธิบายลักษณะพลวัตของระบบ แสดงการไหลของการควบคุมจากกิจกรรมหนึ่งไปยังอีกกิจกรรมหนึ่ง

\subsection{กระบวนการจองและการตรวจสอบเวลาซ้อนทับ (Booking \& Time Conflict Detection Flow)}
กระบวนการนี้แสดงตรรกะที่ใช้เพื่อให้มั่นใจในความถูกต้องของข้อมูลระหว่างการจอง ระบบต้องตรวจสอบว่าทั้งเทรนเนอร์และอุปกรณ์ที่เลือกไม่ถูกใช้งานในช่วงเวลาที่ร้องขอ

\begin{figure}[h]
\centering
\begin{lstlisting}[language=java, caption={Mermaid Code for Booking Flow}, frame=single]
graph TD
    A[Start: Member requests booking] --> B(Select Date & Time)
    B --> C{Calculate Time Interval}
    C --> D[Fetch Existing Sessions]
    D --> E{Is Trainer Available?}
    E -- No --> F[Disable Trainer Option]
    E -- Yes --> G[Enable Trainer Option]
    D --> H{Is Equipment Available?}
    H -- No --> I[Disable Equipment Option]
    H -- Yes --> J[Enable Equipment Option]
    F --> K(User Selects Resources)
    G --> K
    I --> K
    J --> K
    K --> L{Validation Passed?}
    L -- No --> M[Show Error]
    L -- Yes --> N[Save to Database]
    N --> O[End: Booking Confirmed]
\end{lstlisting}
\caption{โค้ด Mermaid แสดงกระบวนการจองและการตรวจสอบความขัดแย้ง}
\end{figure}

\textbf{คำอธิบายอย่างละเอียด:}
\begin{enumerate}
    \item \textbf{การเริ่มต้น}: กระบวนการเริ่มเมื่อสมาชิกเปิดหน้าต่างการจอง
    \item \textbf{การเลือกเวลา}: ผู้ใช้ระบุวันที่ เวลาเริ่ม และระยะเวลา ระบบคำนวณ `endTime` ทันที
    \item \textbf{การสอบถามความขัดแย้ง}: Backend (หรือ Logic ฝั่ง Frontend ผ่าน API) ดึงข้อมูลเซสชันทั้งหมดที่ทับซ้อนกับ `[startTime, endTime]`
    \item \textbf{การกรองทรัพยากร}:
    \begin{itemize}
        \item ระบบตรวจสอบรายการเทรนเนอร์ หาก ID เทรนเนอร์ปรากฏในเซสชันที่ทับซ้อน จะถูกระบุว่า "ไม่ว่าง"
        \item ในทำนองเดียวกัน ID อุปกรณ์ที่พบในเซสชันที่ทับซ้อนจะถูกระบุว่า "ไม่ว่าง"
    \end{itemize}
    \item \textbf{การอัปเดต UI}: Frontend แสดงผล Dropdown ตัวเลือกที่ไม่ว่างจะถูกทำเป็นสีเทา (Disabled)
    \item \textbf{การส่งข้อมูล}: ผู้ใช้เลือกจากตัวเลือกที่ถูกต้องและกดส่ง
    \item \textbf{การตรวจสอบขั้นสุดท้าย}: Server ทำการตรวจสอบครั้งสุดท้ายก่อนบันทึกเพื่อป้องกัน Race Condition หากผ่าน จะบันทึกข้อมูล
\end{enumerate}

\subsection{กระบวนการจัดการข้อมูลของผู้ดูแล (Data Management Flow)}
กระบวนการนี้แสดงวิธีการที่ผู้ดูแลระบบจัดการข้อมูลในระบบ (เช่น เพิ่มเทรนเนอร์ใหม่)

\begin{figure}[h]
\centering
\begin{lstlisting}[language=java, caption={Mermaid Code for Data Management}, frame=single]
graph TD
    A[Start: Admin Dashboard] --> B[Select Entity Type e.g. Trainer]
    B --> C[View List]
    C --> D{Choose Action}
    D -- Add New --> E[Fill Input Form]
    D -- Edit --> F[Modify Existing Data]
    D -- Delete --> G[Confirm Deletion]
    E --> H{Validate Input?}
    F --> H
    H -- Invalid --> I[Show Validation Error]
    H -- Valid --> J[Send API Request]
    G --> J
    J --> K[Database Operation]
    K --> L[Log System Action]
    L --> M[Update UI List]
    M --> N[End]
\end{lstlisting}
\caption{โค้ด Mermaid แสดงกระบวนการจัดการข้อมูลของผู้ดูแลระบบ}
\end{figure}

\textbf{คำอธิบายอย่างละเอียด:}
\begin{enumerate}
    \item \textbf{การเลือก}: Admin เลือกหมวดหมู่ (Trainers, Customers, หรือ Equipment)
    \item \textbf{การเริ่มการกระทำ}: Admin เลือกที่จะ เพิ่ม, แก้ไข, หรือ ลบ
    \item \textbf{การตรวจสอบ}: สำหรับการเพิ่ม/แก้ไข ระบบตรวจสอบข้อกำหนด (เช่น ชื่อห้ามว่าง, อีเมลต้องไม่ซ้ำ)
    \item \textbf{การดำเนินการ}: ส่งคำขอไปยัง API Server ดำเนินการ CRUD กับฐานข้อมูล PostgreSQL
    \item \textbf{การบันทึก}: เมื่อสำเร็จ ฟังก์ชัน `logAction()` จะบันทึกเหตุการณ์
    \item \textbf{การตอบสนอง}: รายการใน UI จะถูกรีเฟรชเพื่อแสดงการเปลี่ยนแปลงทันที
\end{enumerate}

\newpage
\section{แผนภาพลำดับ (Sequence Diagrams)}
แผนภาพลำดับแสดงรายละเอียดการโต้ตอบทางเทคนิคระหว่าง Frontend (View), Backend (Controller) และ Database (Model) ตามลำดับเวลา

\subsection{กระบวนการเข้าสู่ระบบของผู้ใช้ (The User Login Process)}
แผนภาพนี้แสดงขั้นตอนการยืนยันตัวตนตั้งแต่ HTTP Request แรกจนถึงการสร้าง Session ฝั่งไคลเอนต์

\begin{figure}[h]
\centering
\begin{lstlisting}[language=java, caption={Mermaid Code for Login Process}, frame=single]
sequenceDiagram
    participant User
    participant Browser (Frontend)
    participant Server (Express.js)
    participant Database (PostgreSQL)

    User->>Browser: Enter Email & Phone
    User->>Browser: Click "Sign In"
    Browser->>Server: POST /api/login {email, phone}
    activate Server
    Server->>Database: SELECT * FROM Customers WHERE email = ?
    activate Database
    Database-->>Server: Return User Record
    deactivate Database
    
    alt User Found
        Server->>Server: Validate Phone Number
        Server-->>Browser: HTTP 200 OK {user, role}
        Browser->>Browser: localStorage.setItem('user', data)
        Browser->>User: Redirect to Dashboard
    else User Not Found
        Server-->>Browser: HTTP 401 Unauthorized
        Browser->>User: Show Error Message
    end
    deactivate Server
\end{lstlisting}
\caption{โค้ด Mermaid แสดงลำดับการเข้าสู่ระบบ}
\end{figure}

\textbf{คำอธิบายลำดับเหตุการณ์:}
\begin{enumerate}
    \item \textbf{ผู้ใช้} กรอกข้อมูลในฝั่งไคลเอนต์
    \item \textbf{เบราว์เซอร์} ส่งคำขอ `fetch` แบบ Asynchronous (POST) ไปยัง `/api/login`
    \item \textbf{เซิร์ฟเวอร์} รับคำขอและใช้ Sequelize ค้นหาใน \textbf{ฐานข้อมูล}
    \item \textbf{ฐานข้อมูล} ส่งคืนระเบียนผู้ใช้ที่ตรงกับอีเมล
    \item \textbf{เซิร์ฟเวอร์} เปรียบเทียบเบอร์โทรศัพท์ที่ระบุกับที่เก็บไว้
    \item หากถูกต้อง เซิร์ฟเวอร์ตอบกลับด้วย JSON Object ที่มีโปรไฟล์และบทบาทของผู้ใช้
    \item \textbf{เบราว์เซอร์} เก็บข้อมูล Session นี้ใน `localStorage` เพื่อคงสถานะและนำทางผู้ใช้ไปต่อ
\end{enumerate}

\subsection{กระบวนการจอง (The Booking Process)}
แผนภาพนี้แสดงขั้นตอนตั้งแต่ต้นจนจบของการสร้างเซสชันการฝึกซ้อมใหม่

\begin{figure}[h]
\centering
\begin{lstlisting}[language=java, caption={Mermaid Code for Booking Process}, frame=single]
sequenceDiagram
    participant Member
    participant UI (Booking Modal)
    participant API (Controller)
    participant DB (TrainingSessions)

    Member->>UI: Select Date, Time, Duration
    UI->>API: GET /api/sessions?date=...
    activate API
    API->>DB: SELECT * FROM TrainingSessions WHERE date = ...
    activate DB
    DB-->>API: Return List of Sessions
    deactivate DB
    API-->>UI: Return Occupied Slots
    deactivate API
    
    UI->>UI: Filter Available Trainers/Equipment
    Member->>UI: Select Trainer & Equipment
    Member->>UI: Confirm Booking
    
    UI->>API: POST /api/sessions {userId, trainerId, ...}
    activate API
    API->>DB: INSERT INTO TrainingSessions VALUES (...)
    activate DB
    DB-->>API: Success (New ID)
    deactivate DB
    API-->>UI: HTTP 201 Created
    deactivate API
    
    UI->>Member: Show Success Alert
    UI->>UI: Refresh Schedule Grid
\end{lstlisting}
\caption{โค้ด Mermaid แสดงลำดับการจอง}
\end{figure}

\textbf{คำอธิบายลำดับเหตุการณ์:}
\begin{enumerate}
    \item \textbf{การตรวจสอบความพร้อม}: เมื่อสมาชิกเลือกเวลา \textbf{UI} ร้องขอข้อมูลเซสชันที่มีอยู่จาก \textbf{API} เพื่อตรวจสอบความขัดแย้ง
    \item \textbf{การกรอง}: ตรรกะของ \textbf{UI} ใช้ผลลัพธ์จาก API เพื่อปิดการใช้งานตัวเลือกที่ไม่ว่าง
    \item \textbf{การส่งข้อมูล}: สมาชิกกดส่งแบบฟอร์มการจอง
    \item \textbf{การบันทึก}: \textbf{API} รับคำขอ POST และใช้เมธอด `create()` ของ Sequelize เพื่อเพิ่มแถวใหม่ใน \textbf{DB}
    \item \textbf{การยืนยัน}: เมื่อบันทึกสำเร็จ ฐานข้อมูลยืนยันธุรกรรม API ส่งสถานะ 201 กลับไปยังไคลเอนต์ เพื่อแสดงแจ้งเตือนความสำเร็จ
\end{enumerate}
